\documentclass[conference]{IEEEtran}
\usepackage[hidelinks]{hyperref}
\hypersetup{pdfauthor={}, pdftitle={COALESCE}, pdfsubject={}, pdfkeywords={}}
\usepackage{graphicx}
\usepackage{amsmath}
\usepackage{adjustbox}
\usepackage{cite}
\usepackage{algorithm}

\usepackage{booktabs} % For better tables
\usepackage{microtype} % Better typography
\usepackage{xcolor}   % For colored text if needed
\usepackage{float}
\usepackage{listings}
\usepackage{algpseudocode}

\usepackage{tikz}
\usetikzlibrary{decorations.pathreplacing}

\definecolor{codegray}{rgb}{0.5,0.5,0.5}
\definecolor{backcolor}{rgb}{0.95,0.95,0.95}

\title{COALESCE: Coherence-Aware Perceptron Cache Replacement for Multithreaded Workloads}

\author{}

\begin{document}
\raggedbottom
\maketitle

\begin{abstract}
Learning-based last-level cache (LLC) replacement policies such as Hawkeye and Mockingjay define the state of the art, but they are evaluated almost exclusively on multi-programmed workloads---independent programs that share cache capacity without sharing data. We show that this evaluation regime hides a blind spot: under genuine multithreaded sharing, the state of the art collapses. Mockingjay, the strongest policy across one hundred multi-programmed mixes in its original evaluation, finishes last on three of the four multithreaded workloads in our study---up to 48\% behind SRRIP---and Hawkeye falls behind LRU on two of the four. We present COALESCE, a coherence-aware perceptron replacement policy whose features---program counter, MESI coherence state, and sharer count---are designed for coherent workloads. To evaluate it, we extend the ChampSim simulator, whose default per-CPU address spaces silently preclude all inter-core data sharing, with a shared-virtual-memory overlay that exposes coherence behaviour to the replacement policy. Across the four-workload PARSEC and SPLASH-3 suite at 4 cores, COALESCE achieves geometric-mean speedups of 9.7\% over LRU, 1.5\% over Hawkeye, and 13.9\% over Mockingjay, winning outright on canneal, the write-intensive irregular workload in our suite (5\% over Hawkeye and 30\% over Mockingjay at 8 cores). A mechanism decomposition shows the gain comes from coherence-state-aware retention rather than the perceptron itself: a bias-free perceptron matches LRU, while an explicit MODIFIED-state bias supplies the advantage and the sharer feature adds a stable 7\% under genuine cross-thread sharing. We characterise the workload properties---LLC pressure, access irregularity, and write fraction---that predict when coherence-aware replacement helps, and show that when they are absent, simple heuristics remain unbeaten.
\end{abstract}

\section{Introduction}

Modern multi-core processors share their last-level cache (LLC) among all cores, and the policy that decides which block to evict on a miss has a first-order effect on system throughput~\cite{hennessy2017quantitative}. A long line of work has progressively replaced static heuristics such as Least Recently Used (LRU) and Re-Reference Interval Prediction (RRIP)~\cite{jaleel2010rrip} with predictive, learning-based policies: SHiP learns reuse behaviour per program-counter signature~\cite{wu2011ship}, Hawkeye learns to imitate Belady's provably optimal MIN algorithm~\cite{belady1966study,jain2016hawkeye}, and Mockingjay refines this mimicry with fine-grained reuse-distance prediction, approaching the performance of MIN itself on its evaluation suite~\cite{shah2022mockingjay}.

These policies share an evaluation methodology as well as a lineage. They are validated on \emph{multi-programmed} workloads: collections of independent, single-threaded programs (typically drawn from SPEC CPU and GAP) co-scheduled on a multi-core simulator. In such workloads the cores \emph{compete} for LLC capacity but never \emph{share} data: no line is ever read by one core and written by another, no coherence invalidation ever removes a block, and the coherence state of every LLC line is trivial. Multithreaded applications---the very workloads for which shared LLCs are designed---obey different rules: threads communicate through shared lines, writes by one core invalidate copies held by others, and a block's coherence role carries information about its future reuse that no program-counter signature can see.

This paper asks a simple question: \emph{do state-of-the-art learned replacement policies, designed and tuned on multi-programmed workloads, transfer to genuine multithreaded sharing?} Answering it required removing an infrastructure obstacle. ChampSim~\cite{champsim2022}, the \emph{de facto} standard simulator of the cache-replacement community, assigns each simulated CPU a \emph{private} physical address space: two cores accessing the same virtual address receive different physical pages by construction. Every multi-core evaluation built on default ChampSim---including those of the policies above---is therefore multi-programmed, even when the input traces come from a multithreaded program. We extend ChampSim's virtual memory subsystem with a \emph{shared-VMEM overlay} that maps identical virtual addresses from participating cores to the same physical page, restoring genuine inter-core sharing, coherence invalidations, and sharer tracking at the LLC (Section~\ref{sec:methodology}).

With this infrastructure we evaluate seven policies---LRU, SRRIP, DRRIP, SHiP, Hawkeye, Mockingjay, and our proposed COALESCE---on four multithreaded benchmarks from PARSEC~\cite{bienia2008parsec} and SPLASH-3~\cite{sakalis2016splash3} at 4, 8, and 16 cores. The results expose a clear blind spot. Mockingjay, the strongest policy in the multi-programmed world, finishes \emph{last} on three of the four workloads we test, trailing SRRIP by up to 48\%; on the fourth it is second-to-last among the baselines. Hawkeye falls behind even LRU on two workloads. On these workloads, the implicit assumption that multi-programmed performance predicts multithreaded performance does not hold.

We then present COALESCE (Coherence-Observant Adaptive Learning for System-wide Cache Efficiency), a hashed-perceptron replacement policy~\cite{jimenez2001perceptron,jimenez2017multiperspective} whose feature set is chosen for coherent workloads: the program counter of the instruction that filled a block, the block's MESI coherence state, and its current sharer count. COALESCE wins outright on canneal, the most memory-intensive workload in our suite: 5.0\% over Hawkeye and 30.5\% over Mockingjay at 8 cores, and 47\% over LRU at 4 cores. On both workloads where capacity pressure and write traffic coincide (canneal and ocean), it is the strongest of the learning-based policies. Across the full suite---including the workloads where it loses---COALESCE achieves geometric-mean speedups of 9.7\% over LRU and 1.5\% over Hawkeye at 4 cores.

This paper makes four contributions:
\begin{itemize}
  \item \textbf{A methodology contribution}: a shared-virtual-memory extension to ChampSim that enables, to our knowledge, the first evaluation of learned LLC replacement under genuine multithreaded data sharing (Section~\ref{sec:methodology}).
  \item \textbf{An evaluation finding}: state-of-the-art learned policies do not transfer to coherent workloads. Mockingjay places last on three of the four workloads (second-to-last among baselines on the fourth); Hawkeye loses to LRU on two (Section~\ref{sec:results}).
  \item \textbf{A policy}: COALESCE, a coherence-aware perceptron replacement policy that is the strongest learning-based policy on the capacity-pressured, write-active workloads in our suite, with a policy storage budget of ${\sim}44$\,KB for a 2\,MB LLC (Section~\ref{sec:design}).
  \item \textbf{A characterisation}: a feature ablation and workload analysis showing that the sharer-count feature contributes a stable 7\% exactly when cross-thread sharing is genuine, and identifying the three workload properties---LLC pressure, access irregularity, and write fraction---that predict when coherence-aware replacement helps (Sections~\ref{sec:ablation} and~\ref{sec:results}).
\end{itemize}

\section{Background and Motivation}
\label{sec:background}

\subsection{Learned cache replacement and its evaluation}

Belady's MIN algorithm~\cite{belady1966study} evicts the block whose next use is furthest in the future and is provably optimal, but requires knowledge of the future. Modern learned policies approximate it from history. SHiP~\cite{wu2011ship} associates re-reference behaviour with program-counter (PC) signatures. Hawkeye~\cite{jain2016hawkeye} reconstructs what MIN \emph{would have done} on a sampled history window (OPTgen) and trains a PC-indexed predictor to label load instructions cache-friendly or cache-averse. Mockingjay~\cite{shah2022mockingjay} replaces Hawkeye's binary labels with predicted reuse distances, tracked per PC in a sampled cache, and reports performance within 0.3\% of MIN's improvement bound on a single core. Glider~\cite{shi2019glider} applies offline deep learning to discover better features. Perceptron-based reuse prediction~\cite{jimenez2001perceptron,jimenez2017multiperspective} hashes program features into weight tables whose summed response drives the eviction decision.

All of these policies are evaluated on single-threaded traces (SPEC CPU 2006/2017, GAP, CVP) run either alone or as multi-programmed multi-core mixes. We are not aware of any published evaluation of these policies under multithreaded data sharing.

\subsection{What sharing changes}

In a multithreaded workload, the LLC is not merely contended capacity; it is the communication fabric between threads. This changes replacement in three ways. First, \emph{coherence state carries reuse information}: a block in MODIFIED state~\cite{sorin2011coherence} was recently written and---in producer-consumer and migratory sharing patterns---is disproportionately likely to be read soon by another core. Evicting it penalises the future reader and forces a writeback. Second, \emph{sharer count carries reuse information}: a line cached by several cores' private levels is part of the actively shared working set; evicting its LLC copy forces all sharers to refetch from memory. Third, \emph{invalidations break recency heuristics}: a write by one core removes copies from other cores' private caches, so LLC recency no longer reflects the line's system-wide liveness. None of these signals exist in a multi-programmed workload, where every line has exactly one ever-owner, no line is ever invalidated, and state never carries cross-core meaning.

\subsection{The simulator obstacle}
\label{sec:obstacle}

ChampSim~\cite{champsim2022} keys its virtual-to-physical mapping on the pair (CPU, virtual address): two cores touching the same virtual address are assigned \emph{different} physical pages. This is convenient for multi-programmed studies---identical SPEC traces can be replicated across cores without aliasing---but it makes genuine sharing impossible: even traces captured from one multithreaded process are silently privatised, every sharer count is 1, and the number of coherence invalidations at the LLC is exactly zero on every run. We verified this empirically on our workloads; it motivates the overlay described next.

\section{The COALESCE Policy}
\label{sec:design}

COALESCE is a hashed-perceptron replacement policy in the lineage of multiperspective reuse prediction~\cite{jimenez2017multiperspective}, with a feature set chosen for coherent workloads. It comprises two perceptron weight tables, a per-sampled-set ghost filter for eviction feedback, and an explicit coherence bias; Table~\ref{tab:storage} itemises the structures.

\subsection{Prediction structures}

Two independent weight tables of 2{,}048 8-bit saturating entries each ($[-128, +127]$) are indexed by orthogonal hashes of the candidate block's metadata:
\begin{equation}
h_0 = \mathrm{hash}\big(\mathrm{PC} \oplus \mathrm{0x9e3779b9},\ \mathrm{state}\big) \bmod 2048,
\end{equation}
\begin{equation}
h_1 = \mathrm{hash}\big(\mathrm{PC} \oplus \mathrm{0x85ebca6b},\ \mathrm{sharers}\big) \bmod 2048,
\end{equation}
where PC is the program counter of the instruction that most recently filled or touched the block, \emph{state} is the block's MESI state, and \emph{sharers} is the population count of its sharer mask. The raw vote for a block is the sum $w_0[h_0] + w_1[h_1]$: large positive votes mean ``keep'', large negative votes mean ``safe to evict''.

\subsection{Coherence bias}

On top of the learned vote, COALESCE applies an explicit coherence bias, but only to blocks the perceptron already considers worth keeping (raw vote $> 0$):
\begin{equation}
\mathrm{vote} \mathrel{+}=
\begin{cases}
+40 & \text{if state} = \text{MODIFIED},\\
+20 \times \mathrm{sharers} & \text{if sharers} \geq 2.
\end{cases}
\label{eq:bias}
\end{equation}
The gating prevents the bias from rescuing blocks the perceptron has confidently learned are dead. The magnitudes were set empirically from a sweep on 8-core canneal (Section~\ref{sec:ablation}); within a broad neighbourhood of the chosen $(40, 20)$ the policy is insensitive to the exact values, but the bias as a whole is essential---without it the policy degrades to the level of LRU. The victim is the way with the \emph{lowest} final vote (Algorithm~\ref{alg:victim}).

\begin{algorithm}[tbp]
\small
\caption{Victim Selection in COALESCE}
\label{alg:victim}
\begin{algorithmic}[1]
\Procedure{FindVictim}{$set$}
    \State $victim \gets 0$, $minVote \gets +\infty$
    \For{$way \gets 0$ to $NUM\_WAY - 1$}
        \If{$\neg\, block[way].valid$} \Return $way$ \EndIf
        \State $v \gets w_0[h_0(\mathrm{PC}, \mathrm{state})] + w_1[h_1(\mathrm{PC}, \mathrm{sharers})]$
        \If{$v > 0$}
            \If{$state = \mathrm{MODIFIED}$} $v \gets v + 40$ \EndIf
            \If{$sharers \geq 2$} $v \gets v + 20 \times sharers$ \EndIf
        \EndIf
        \If{$v < minVote$} $minVote \gets v$;\ $victim \gets way$ \EndIf
    \EndFor
    \If{$set$ is sampled} record victim in ghost filter \EndIf
    \State \Return $victim$
\EndProcedure
\end{algorithmic}
\end{algorithm}

\subsection{Training and the ghost-rescue mechanism}

Training is confined to sampled sets (every 32nd set, 3.125\% of the cache) to bound metadata cost, following the set-sampling tradition of DRRIP and Hawkeye~\cite{jaleel2010rrip,jain2016hawkeye}. On a hit in a sampled set, the weights at $h_0$ and $h_1$ are incremented (reward) when the current vote mispredicts or has magnitude below a confidence threshold $\theta = 35$, a training rule adapted from perceptron branch prediction~\cite{jimenez2001perceptron}; $\theta$ was fixed a priori and not tuned per workload.

Training on misses occurs through a \emph{ghost rescue} path. Each sampled set maintains a 1{,}024-bit Bloom filter (3 hash functions, reset every 150 insertions)\footnote{With $k{=}3$ hashes, $m{=}1024$ bits, and $n{=}150$ insertions between resets, the worst-case false-positive rate is $(1-e^{-kn/m})^k \approx 4.5\%$; the periodic reset bounds both the false-positive rate and entry staleness.} and a 128-entry direct-mapped ghost table of recently evicted blocks; each 32-bit ghost entry packs a 12-bit PC signature, 14-bit partial tag, 3-bit sharer count, 2-bit MESI state, and a valid bit. When a miss matches a ghost entry---evidence that the policy evicted a block that was still live---the corresponding weights receive a boosted reward ($5\times$ repetition), rapidly rehabilitating the mispredicted context.

\subsection{Storage cost}
\label{sec:storage}

Table~\ref{tab:storage} itemises the policy storage. MESI state and the sharer mask are maintained by the coherence substrate of any real multi-core (directory or snoop filter) and are read, not duplicated, by COALESCE; we therefore count only policy-private structures, as is conventional~\cite{jain2016hawkeye,shah2022mockingjay}. The total is ${\sim}44$\,KB for a 2\,MB LLC ($\sim$2.1\% of the data array). For comparison, Mockingjay reports 31.91\,KB on one core~\cite{shah2022mockingjay} and Hawkeye fits the 32\,KB replacement-championship budget~\cite{jain2016hawkeye}; COALESCE's overhead is modestly larger, dominated by the ghost tables, and the ablation of Section~\ref{sec:ablation} indicates the policy's advantage is carried by its feature set rather than by this extra state.

\begin{table}[htbp]
    \centering
    \renewcommand{\arraystretch}{1.2}
    \caption{COALESCE policy storage for a 2\,MB, 2048-set LLC}
    \label{tab:storage}
    \resizebox{0.95\linewidth}{!}{%
    \begin{tabular}{lll}
        \toprule
        \textbf{Structure} & \textbf{Sizing} & \textbf{Cost} \\
        \midrule
        Perceptron tables & $2 \times 2048 \times 8$\,bits & 4\,KB \\
        Bloom filters & $64$ sampled sets $\times\ 1024$\,bits & 8\,KB \\
        Ghost tables & $64 \times 128$ entries $\times\ 32$\,bits & 32\,KB \\
        \midrule
        \textbf{Total} & & \textbf{44\,KB} \\
        \bottomrule
    \end{tabular}
    }
\end{table}

\section{Experimental Methodology}
\label{sec:methodology}

\subsection{Shared-VMEM overlay}

We extend ChampSim's \texttt{VirtualMemory} with a configurable set of \emph{sharing CPUs}. For CPUs in this set, the virtual-to-physical mapping is keyed on the virtual address alone, so identical virtual addresses from different cores resolve to the same physical page---restoring the semantics of threads within one process, i.e., the sharing behaviour that PARSEC and SPLASH-3 were designed to exercise~\cite{bienia2008parsec,sakalis2016splash3}. Fills that hit a page mapped by another core are tracked (\emph{aliased fills}), the LLC maintains a per-line sharer mask and MESI state, and writes to lines with remote sharers generate coherence invalidations, all of which are exposed to the replacement policy and reported in statistics. The overlay is independent of any policy. We release the overlay, all seven policy implementations, and every run configuration as an artifact\footnote{Anonymised for review: \url{https://anonymous.4open.science/r/COALESCE-Hashed-Perceptron-Cache-Replacement-Policy-1E6D/}}; all results in this paper are reproducible from the supplied configurations, and ChampSim's determinism makes the reported cycle counts exact for a given configuration, which is why tables report unrounded values. All experiments in this paper, for all policies, run under the overlay; Section~\ref{sec:threats} discusses the default-VMEM (private) configuration.

We validate the overlay by confirming that its coherence statistics track each workload's write semantics. Across every workload the overlay records nonzero cross-CPU aliased fills, confirming it is active; meanwhile LLC coherence invalidations are \emph{exactly zero} on fluidanimate (100\% LLC reads) at both 4 and 8 cores, and scale with core count on write-heavy canneal---0.11\,M, 6.9\,M, and 23.9\,M invalidations at 4, 8, and 16 cores respectively. An overlay fabricating sharing could not produce exactly-zero invalidations on a read-only workload alongside core-count-proportional invalidations on a write-heavy one; together these establish that the recorded coherence activity is genuine.

\subsection{System configuration}

Table~\ref{tab:simconfig} lists the simulated system. The LLC is shared by all cores; private L1/L2 caches and TLBs use ChampSim defaults. No prefetcher is used at any level, isolating the replacement policy's effect, consistent with the no-prefetcher configurations of prior replacement studies~\cite{jain2016hawkeye,shah2022mockingjay}. The 50\,M-instruction-per-core warmup (0.4--1.6\,B committed instructions per run) cycles the 32K-line LLC many times over, and per-core heartbeat IPC stabilises well before the measurement window in every run.

\begin{table}[htbp]
    \centering
    \caption{Simulation Configuration}
    \label{tab:simconfig}
    \resizebox{0.95\linewidth}{!}{%
    \begin{tabular}{ll}
        \toprule
        \textbf{Parameter} & \textbf{Value} \\
        \midrule
        Cores & 4, 8, 16 (out-of-order, ChampSim default core) \\
        Shared LLC & 2\,MB, 2048 sets $\times$ 16 ways, 64\,B lines \\
        Private L1I/L1D/L2 & ChampSim defaults per core \\
        DRAM & DDR4-3200, $t_{RP}=t_{RCD}=t_{CAS}=24$ (ChampSim default) \\
        Virtual memory & shared-VMEM overlay across all cores \\
        Prefetcher & none \\
        Warmup / simulation & 50\,M / 100\,M instructions per core \\
        Replacement policies & LRU, SRRIP, DRRIP, SHiP, Hawkeye, \\
        & Mockingjay, COALESCE (+ ablation) \\
        \bottomrule
    \end{tabular}
    }
\end{table}

\subsection{Workloads}

We use four multithreaded benchmarks chosen to span distinct LLC behaviours (Table~\ref{tab:workloads}): \texttt{canneal} (PARSEC; simulated-annealing netlist routing---pointer-chasing, write-heavy, far exceeding LLC capacity), \texttt{ocean\_cp} (SPLASH-3; multigrid PDE solver---regular grid sweeps, mixed read/write, capacity-pressured), \texttt{fluidanimate} (PARSEC; particle simulation whose LLC traffic is 100\% reads), and \texttt{barnes} (SPLASH-3; N-body octree whose working set at the standard $n{=}16384$ input fits within the LLC). Traces are captured per thread with a Pin-based tracer~\cite{luk2005pin} and replayed with one thread per core; at 8 and 16 cores, thread traces are replicated in the same pattern for all policies, so all comparisons are like-for-like. Within each benchmark and core count, every policy executes an identical instruction volume.

\begin{table}[htbp]
    \centering
    \renewcommand{\arraystretch}{1.2}
    \caption{Workload characterisation at the shared LLC (8 cores)}
    \label{tab:workloads}
    \resizebox{\linewidth}{!}{%
    \begin{tabular}{lllll}
        \toprule
        \textbf{Property} & \texttt{canneal} & \texttt{ocean} & \texttt{fluidanim.} & \texttt{barnes} \\
        \midrule
        Access pattern & irregular & regular & regular & irregular \\
        LLC write fraction & ${\sim}89\%$ & ${\sim}40\%$ & $0\%$ & mixed \\
        Working set vs.\ LLC & $\gg$ & $\gg$ & $>$ & $<$ \\
        LLC MPKI (range) & 13--31 & 37--146 & $<0.3$ & ${\sim}1$ \\
        LLC invalidations & 6.9\,M & 0.72\,M & 0 & 0.03\,M \\
        Lines with $\geq$2 sharers & 25.6\% & high & ${\sim}0$ & low \\
        \bottomrule
    \end{tabular}
    }
\end{table}

\subsection{Baselines and metrics}

LRU, SRRIP, DRRIP, and SHiP use ChampSim's standard implementations. We implemented Hawkeye (OPTgen with sampled history plus a PC-indexed predictor)~\cite{jain2016hawkeye} and Mockingjay (sampled reuse-distance prediction with estimated time remaining)~\cite{shah2022mockingjay} as ChampSim modules following the original papers; Table~\ref{tab:baselines} lists the exact configurations, which follow each paper's published design so that no policy is starved of metadata.

\begin{table}[htbp]
    \centering
    \renewcommand{\arraystretch}{1.2}
    \caption{Reimplemented baseline configurations}
    \label{tab:baselines}
    \resizebox{\linewidth}{!}{%
    \begin{tabular}{ll}
        \toprule
        \textbf{Hawkeye}~\cite{jain2016hawkeye} & 8K-entry PC predictor, 3-bit counters; \\
        & 64 sampled sets $\times$ 8-way history; \\
        & OPTgen window $8\times$ associativity; 3-bit RRPV \\
        \midrule
        \textbf{Mockingjay}~\cite{shah2022mockingjay} & 16K-entry RDP per CPU, 6-bit values; \\
        & 64 sampled sets $\times$ 8-way sampled cache; \\
        & 6-bit ETR counters, ageing quantum 16 accesses \\
        \bottomrule
    \end{tabular}
    }
\end{table}

Our primary metric is \emph{total execution cycles}: the cycle count at which the last core completes its 100\,M-instruction region (lower is better). The weighted-speedup metric conventional in multi-programmed studies~\cite{shah2022mockingjay} is not meaningful here: it normalises each program's shared-cache IPC by its isolated-execution IPC, a construction that presumes independent programs, whereas our threads cooperate within a single program and their progress is mutually dependent. For the same reason we report no Belady's MIN upper bound---as the Mockingjay authors themselves note, MIN is not well defined for multi-core systems~\cite{shah2022mockingjay}, and under inter-thread timing dependence the offline-optimal eviction sequence is not unique. Per-core IPC is therefore reported only for bottleneck analysis, and speedups are reported as cycle-count ratios summarised by geometric mean. We additionally report LLC misses per kilo-instruction (MPKI) as a secondary metric, normalised to the instructions of the \emph{LLC-active} cores: cores whose replicated traces fit within their private caches generate no LLC traffic and are excluded from the denominator. The normalisation is identical for every policy, so all comparisons are unaffected. Note that on multithreaded workloads MPKI and execution time can \emph{invert}---Section~\ref{sec:ocean} shows LRU attaining the lowest MPKI on ocean while finishing 6--7\% slower than SHiP---because misses differ in criticality and writeback cost when threads contend on a coherent LLC. This inversion is itself a reason total cycles must be the primary metric in this setting.

\section{Results}
\label{sec:results}

\subsection{canneal: where coherence-aware learning wins}
\label{sec:canneal}

Figure~\ref{fig:canneal} and Table~\ref{tab:canneal} report total cycles and LLC MPKI on canneal. At 4 cores COALESCE finishes in 267.3\,M cycles---7.3\% ahead of Hawkeye, the best baseline, and 47.1\% ahead of LRU---and the mechanism is visible in raw misses: 12.9 MPKI against 24.1--30.9 for every baseline, roughly \emph{half} the misses of the best competitor. At 8 cores COALESCE finishes first at 301.9\,M cycles; its 0.4\% margin over SRRIP is within plausible run-to-run variation and we treat the two as tied pending the variance study of Section~\ref{sec:variance}, but its margins over SHiP (4.8\%), Hawkeye (5.0\%), and LRU and Mockingjay (30\%) are unambiguous.

canneal is precisely the workload class COALESCE targets: its LLC traffic is ${\sim}89\%$ writes (RFO + writeback)\footnote{Measured from CPU~0's LLC demand and writeback stream at 8 cores: 88.8\% under every policy tested---the access mix is policy-invariant to within 0.1\%.}, its pointer-chasing access pattern defeats recency and RRIP-style position heuristics, and its working set vastly exceeds the 2\,MB LLC, so every eviction decision matters. Under the shared overlay its coherence machinery is highly active (6.9\,M invalidations; one quarter of evicted lines have two or more sharers at 8 cores), giving the MESI and sharer features genuine signal.

\begin{table}[htbp]
    \centering
    \renewcommand{\arraystretch}{1.15}
    \caption{canneal total cycles (lower is better)}
    \label{tab:canneal}
    \resizebox{0.9\linewidth}{!}{%
    \begin{tabular}{lrrrrrr}
        \toprule
        \textbf{Policy} & \textbf{4-core} & vs.\ C. & MPKI & \textbf{8-core} & vs.\ C. & MPKI \\
        \midrule
        \textbf{COALESCE} & \textbf{267{,}251{,}645} & --- & \textbf{12.9} & \textbf{301{,}945{,}789} & --- & 13.3 \\
        Hawkeye & 286{,}865{,}528 & $+7.3\%$ & 24.1 & 317{,}172{,}523 & $+5.0\%$ & 19.4 \\
        SRRIP & 287{,}031{,}501 & $+7.4\%$ & 24.3 & 303{,}011{,}899 & $+0.4\%$ & 13.1 \\
        SHiP & 288{,}163{,}881 & $+7.8\%$ & 24.8 & 316{,}393{,}994 & $+4.8\%$ & 15.2 \\
        DRRIP & 294{,}947{,}400 & $+10.4\%$ & 27.6 & 325{,}848{,}199 & $+7.9\%$ & 16.9 \\
        Mockingjay & 311{,}277{,}369 & $+16.5\%$ & 30.8 & 393{,}934{,}080 & $+30.5\%$ & 30.6 \\
        LRU & 392{,}999{,}637 & $+47.1\%$ & 30.9 & 392{,}564{,}736 & $+30.0\%$ & 30.9 \\
        \bottomrule
    \end{tabular}
    }
\end{table}

\begin{figure}[H]
    \centering
    \includegraphics[width=0.95\linewidth]{fig_canneal.png}
    \caption{canneal total execution cycles at 4 and 8 cores (lower is better). COALESCE is first at both scales.}
    \label{fig:canneal}
\end{figure}

The win concentrates exactly where it gates performance. canneal's execution time is set by its two bottleneck cores, and COALESCE serves them best: at 8 cores their IPC is ${\sim}0.332$, a 5.5\% lift over Hawkeye (0.316) and 30\% over Mockingjay and LRU (0.255) on precisely the cores that determine total cycles.

At 16 cores under the overlay, COALESCE completes the same per-core work in 348.2\,M cycles, with bottleneck-core IPC declining from 0.374 at 4 cores to 0.332 at 8 and 0.294 at 16 as per-core effective LLC capacity shrinks by $4\times$ (Figure~\ref{fig:scaling}). We report per-policy comparisons at 4 and 8 cores; full 16-core baselines are left to an extended evaluation, as per-core effective capacity at 16 cores (128\,KB) enters a regime where capacity, rather than the replacement decision, increasingly dominates.

\begin{figure}[tbp]
    \centering
    \includegraphics[width=0.95\linewidth]{fig_scaling.png}
    \caption{COALESCE scaling on canneal: total cycles and bottleneck-core IPC at 4, 8, and 16 cores.}
    \label{fig:scaling}
\end{figure}

\subsection{ocean: regular access favours RRIP heuristics}
\label{sec:ocean}

On ocean (Table~\ref{tab:ocean}) the RRIP family wins: SHiP, SRRIP, and DRRIP finish within 0.4\% of one another, with COALESCE fourth, 3.6\% behind at 4 cores and 2.9\% behind at 8 cores. ocean's multigrid sweeps are regular and stride-predictable---exactly the re-reference geometry that RRIP's insertion heuristics were designed around~\cite{jaleel2010rrip}---and although its coherence traffic is substantial ($6\times$ canneal's invalidation rate per LLC access), regularity, not sharing, dominates the eviction decision.

The MPKI column of Table~\ref{tab:ocean} carries two observations. First, LRU attains the \emph{lowest} miss count at both scales (74.0 and 36.7 MPKI) yet finishes 6--7\% behind SHiP: on a coherent LLC, miss counts and performance can invert, because LRU's recency bias retains recently-written dirty lines whose writebacks retire off the critical path, while evicting lines whose demand refetches stall progress. Second, Mockingjay's degradation is a genuine miss explosion---146.1 MPKI, $1.7\times$ SHiP's---not a latency artefact, which further argues against an implementation-fidelity explanation.

The notable result is the fate of the \emph{learning} policies. COALESCE is the best of them at both scales: it beats Hawkeye by 3.5--4.3\% and LRU by ${\sim}3\%$, while Mockingjay degrades sharply, to 48\% behind SRRIP (Section~\ref{sec:mockingjay}). On coherent workloads, learning does not automatically lose to heuristics---but learning \emph{tuned on the wrong distribution} does.

\begin{table}[htbp]
    \centering
    \renewcommand{\arraystretch}{1.15}
    \caption{ocean total cycles (lower is better)}
    \label{tab:ocean}
    \resizebox{0.9\linewidth}{!}{%
    \begin{tabular}{lrrrrrr}
        \toprule
        \textbf{Policy} & \textbf{4-core} & vs.\ best & MPKI & \textbf{8-core} & vs.\ best & MPKI \\
        \midrule
        SHiP & \textbf{549{,}605{,}858} & --- & 84.8 & 550{,}732{,}443 & $+0.01\%$ & 42.7 \\
        SRRIP & 550{,}297{,}702 & $+0.1\%$ & 85.0 & \textbf{550{,}675{,}830} & --- & 42.6 \\
        DRRIP & 551{,}685{,}991 & $+0.4\%$ & 85.3 & 551{,}564{,}285 & $+0.2\%$ & 42.7 \\
        \textbf{COALESCE} & 569{,}457{,}207 & $+3.6\%$ & 89.9 & 566{,}374{,}131 & $+2.9\%$ & 44.6 \\
        LRU & 587{,}215{,}973 & $+6.8\%$ & \textbf{74.0} & 584{,}499{,}393 & $+6.1\%$ & \textbf{36.7} \\
        Hawkeye & 589{,}923{,}943 & $+7.3\%$ & 94.7 & 591{,}008{,}754 & $+7.3\%$ & 47.6 \\
        Mockingjay & 813{,}767{,}568 & $+48.1\%$ & 146.1 & 814{,}696{,}310 & $+47.9\%$ & 73.2 \\
        \bottomrule
    \end{tabular}
    }
\end{table}

\subsection{fluidanimate and barnes: the boundaries}
\label{sec:boundaries}

The remaining two workloads delimit where coherence-aware replacement cannot help, and we report them as such.

\textbf{fluidanimate}'s LLC traffic is 100\% loads: lines are never written at the LLC, MESI state never leaves EXCLUSIVE, invalidations are exactly zero, and both coherence features of COALESCE are structurally unable to fire. The LLC is also nearly irrelevant in absolute terms---MPKI is below 0.3 under every policy, two orders of magnitude below canneal's. The total policy spread is only ${\sim}5\%$, and COALESCE is last within it (4.9--5.0\% behind DRRIP at both 4 and 8 cores)---paying its prediction overhead for features that carry no signal on this workload.\footnote{The fluidanimate 4-core runs use a 200\,M-warmup/50\,M-simulation window rather than the 50\,M/100\,M used elsewhere; the qualitative result (COALESCE last within a ${\sim}5\%$ spread) is identical to the 8-core run under standard windows and is unaffected.} The consistency of this deficit across core counts confirms it is mechanism, not noise.

\textbf{barnes} at the standard $n{=}16384$ input has a working set ($\sim$1.6\,MB) that fits within the 2\,MB LLC: the LLC observes only ${\sim}500$\,K accesses per core, MPKI stays near 1 under every policy, and evictions are rare, so replacement choices barely matter. All eight policies (the seven of Section~\ref{sec:methodology} plus the no-sharer ablation variant of Section~\ref{sec:ablation}) finish within 3.8\% of one another and LRU---the cheapest---wins; COALESCE trails it by 2.5\%. Low LLC pressure is the third boundary condition: when nothing is evicted, no replacement policy can earn its overhead.

\subsection{The multi-programmed state of the art does not transfer}
\label{sec:mockingjay}

Across all four workloads a consistent pattern emerges for the strongest multi-programmed policies (Figure~\ref{fig:geomean}). Mockingjay---which improves on LRU by 15.2\% averaged over one hundred multi-programmed mixes in its original evaluation~\cite{shah2022mockingjay}---finishes \emph{last on three of the four workloads in this study} and second-to-last among the baselines on the fourth: 30.5\% behind COALESCE on canneal at 8 cores, 48\% behind SRRIP on ocean at both scales, and behind LRU on fluidanimate and barnes. Hawkeye fares better but still loses to LRU on ocean and barnes. We hypothesise distribution shift as the cause. Mockingjay's reuse-distance predictor is trained online from PC-indexed samples that, under multithreaded sharing, conflate accesses from different cores to the same lines. Invalidations also remove blocks at times its estimated-time-remaining model cannot anticipate. Confirming this account requires predictor-accuracy instrumentation on shared versus private lines, which we leave to future work. What the data establishes is the headline: a policy that approaches MIN on the distribution it was designed for can sit below LRU outside it---\emph{on multithreaded workloads, evaluation methodology, not predictor capacity, is a primary bottleneck for learned replacement.}

\subsection{Suite-wide summary}

Figure~\ref{fig:geomean} reports geometric-mean speedups over the same four-workload suite at both scales, losses included. Among the learning-based policies, COALESCE has the highest geomean speedup over LRU at both scales---$+9.7\%$ at 4 cores and $+6.2\%$ at 8 cores, against Hawkeye's $+8.1\%$ and $+5.2\%$, while Mockingjay is net negative ($-3.7\%$ and $-9.2\%$). Stated as COALESCE's margin over each baseline, the suite geomean is $+9.7\%$/$+6.2\%$ over LRU and $+13.9\%$/$+17.0\%$ over Mockingjay at 4/8 cores, with $+1.5\%$/$+1.0\%$ over Hawkeye; against the RRIP family COALESCE sits between $+0.2\%$ and $-2.3\%$. The picture is consistent across scales and fair to state plainly: COALESCE wins decisively inside its target class, concedes modest ground to RRIP heuristics outside it, and is the strongest learning-based policy at both scales.

\begin{figure}[H]
    \centering
    \includegraphics[width=0.95\linewidth]{fig_geomean.png}
    \caption{Geometric-mean speedup of COALESCE over each baseline across the four-workload suite (losses included), at 4 and 8 cores.}
    \label{fig:geomean}
\end{figure}

\subsection{Run-to-run variance}
\label{sec:variance}

ChampSim is deterministic for a fixed configuration, so the relevant variation is across configuration perturbations rather than repeated runs. We quantify it by re-running an 8-core canneal configuration under three virtual-memory randomisation seeds---which perturb physical-page assignment and therefore set-index conflict patterns---for the four policies closest to the headline comparison, holding the workload mix and every other parameter fixed. This is a separate variance experiment on its own fixed mix, not error bars on Table~\ref{tab:canneal}, and the absolute cycle counts therefore differ from that table.

\begin{table}[tbp]
    \centering
    \renewcommand{\arraystretch}{1.15}
    \caption{Seed sensitivity: total cycles (millions) across three VMEM randomisation seeds (8-core canneal, fixed mix)}
    \label{tab:variance}
    \resizebox{0.95\linewidth}{!}{%
    \begin{tabular}{lrrrr}
        \toprule
        \textbf{Policy} & \textbf{min} & \textbf{mean} & \textbf{max} & range/mean \\
        \midrule
        \textbf{COALESCE} & \textbf{268.6} & \textbf{272.9} & \textbf{276.3} & 2.8\% \\
        SRRIP & 283.8 & 286.5 & 290.2 & 2.2\% \\
        SHiP & 294.1 & 295.0 & 296.4 & 0.8\% \\
        Hawkeye & 316.2 & 318.5 & 322.4 & 2.0\% \\
        \bottomrule
    \end{tabular}
    }
\end{table}

Three findings follow (Table~\ref{tab:variance}). First, per-policy variation is small: the range across seeds is at most 2.8\% of the mean (coefficient of variation $\leq 1.2\%$). Second, the policy ranking is identical in every seed---COALESCE $<$ SRRIP $<$ SHiP $<$ Hawkeye---and the bands do not overlap: COALESCE's \emph{slowest} seed (276.3\,M) is still 2.7\% faster than SRRIP's \emph{fastest} (283.8\,M). Third, on this mix COALESCE leads SRRIP by 5.0\% in the mean, well outside either policy's seed band. The 0.4\% COALESCE--SRRIP gap on the headline mix (Section~\ref{sec:canneal}) lies \emph{within} COALESCE's own 2.8\% band, which is why we treat that particular comparison as a tie; the variance study shows the underlying ordering is nonetheless stable, and the 5\% margin over Hawkeye and the ${\sim}30\%$ margins over Mockingjay and LRU exceed the seed band by an order of magnitude.

\section{What Carries the Policy? A Mechanism Decomposition}
\label{sec:ablation}

COALESCE combines two perceptron feature axes (PC$\times$MESI in $h_0$, PC$\times$sharers in $h_1$) with two explicit coherence biases ($+40$ MODIFIED, $+20\times$sharers). We decompose which components actually do the work, with two experiments: a per-axis feature ablation, and a sweep of the bias magnitudes.

\subsection{Per-axis ablation}

We build two ablated variants: \emph{no-sharer} replaces $h_1$'s sharer input with a second PC-only hash and removes the $+20\times$sharers bias (keeping the MESI axis); \emph{no-MESI} symmetrically removes the MESI input from $h_0$ and the $+40$ MODIFIED bias (keeping the sharer axis). Table~\ref{tab:ablation} reports both against full COALESCE; Figure~\ref{fig:ablation} plots the sharer-axis contribution.

The two axes specialise by workload. On canneal the MESI axis carries the policy: removing it costs 20\% (362.0\,M vs.\ 301.9\,M), while removing the sharer axis costs nothing (the cross-thread sharing is shallow, only 1.45\% of evicted lines have $\geq 2$ sharers at 4 cores, so the $+20\times$sharers term is inert). On ocean, whose grid-boundary exchange produces genuine multi-sharer lines, \emph{both} axes contribute: removing either costs 6--8\%. Neither feature is dead weight and neither is universally decisive---each activates where its coherence signal is real (fluidanimate is omitted by construction: its 100\% read traffic leaves both axes unable to fire).

\begin{table}[tbp]
    \centering
    \renewcommand{\arraystretch}{1.15}
    \caption{Per-axis ablation: total cycles for full COALESCE and the two single-axis variants (parenthesised: \% slower than full)}
    \label{tab:ablation}
    \resizebox{\linewidth}{!}{%
    \begin{tabular}{lrrr}
        \toprule
        \textbf{Workload} & \textbf{COALESCE} & \textbf{no-sharer} & \textbf{no-MESI} \\
        \midrule
        canneal 8-core & 301{,}945{,}789 & 300{,}714{,}619 ($-0.4$) & 362{,}027{,}049 (\textbf{+20.0}) \\
        ocean 4-core & 569{,}457{,}207 & 611{,}083{,}533 (+7.3) & 605{,}589{,}821 (+6.3) \\
        ocean 8-core & 566{,}374{,}131 & 607{,}309{,}316 (+7.2) & 609{,}770{,}338 (+7.7) \\
        canneal 4-core & 267{,}251{,}645 & 267{,}247{,}657 ($-0.0$) & --- \\
        barnes 4-core & 160{,}448{,}239 & 160{,}656{,}244 (+0.1) & --- \\
        \bottomrule
    \end{tabular}
    }
\end{table}

\begin{figure}[tbp]
    \centering
    \includegraphics[width=0.95\linewidth]{fig_ablation.png}
    \caption{Contribution of the sharer-count feature (positive = full COALESCE faster than the no-sharer ablation).}
    \label{fig:ablation}
\end{figure}

\subsection{Bias sweep: the coherence bias is the engine}

The ablation removes a feature \emph{and} its bias together. To separate the learned perceptron weights from the explicit biases, we sweep the bias magnitudes $(B_M, B_S)$ on the fixed 8-core canneal configuration of Section~\ref{sec:variance} (full COALESCE: 276.3\,M cycles) while leaving the perceptron untouched. The result is striking: with both biases zero---a fully trained perceptron over PC, MESI, and sharer features, with no explicit coherence prior---COALESCE slows to 391.7\,M cycles, the level of LRU, erasing its entire $\sim$30\% advantage on this configuration. The gain comes from the explicit MODIFIED-retention bias, not from perceptron learning: restoring $B_M{=}40$ alone recovers nearly all of it ($391.7 \to 280.4$\,M), the sharer bias alone recovers little ($391.7 \to 374.2$\,M), and the chosen $(40,20)$ reaches 276.3\,M. Within a broad neighbourhood of the chosen point the policy is insensitive to the exact magnitudes; far from it ($B_M{=}0$) it collapses to the bias-free baseline.

\textbf{Interpretation.} COALESCE's advantage is \emph{coherence-state-aware retention realised as an explicit bias}, not perceptron prediction: the perceptron alone matches LRU on canneal, while the MESI bias supplies the gain and the sharer bias adds a workload-dependent increment under genuine sharing. The perceptron gates the bias---applied only to lines already ranked worth keeping---and carries the ocean case; where the coherence signal cannot fire (fluidanimate) the bias has nothing to act on and COALESCE reverts toward the baselines. We report plainly that the coherence prior, not the learner, is the load-bearing component.

\section{Limitations and Discussion}
\label{sec:threats}

\textbf{Workload coverage.} Four multithreaded benchmarks is small compared to the 30--100 single-threaded traces of multi-programmed studies~\cite{shi2019glider,shah2022mockingjay}; each of our runs, however, is a full multi-core coherent simulation rather than a replicated single-core trace, and the four were chosen to span the characterisation space of Table~\ref{tab:workloads} (write-heavy irregular, write-mixed regular, read-only, and low-pressure). Coherence-stress workloads (e.g., GAP graph kernels run multithreaded, SPLASH-3 water) are future work.

\textbf{Simulation regime.} All results use our shared-VMEM overlay. Under default ChampSim (private per-core address spaces) we observe larger absolute gaps in COALESCE's favour, but regard that regime as a simulator artefact and report only the shared one. We model MESI state and invalidation at the LLC; a full directory protocol with private-cache back-invalidation timing may shift constants but not, we believe, the structural findings.

\textbf{Input and cache scale.} barnes uses the standard $n{=}16384$ input (working set fits the LLC; larger inputs hit a static array limit in the reference binary), and all experiments use a single 2\,MB LLC. A working-set/cache-size sweep---which the characterisation of Table~\ref{tab:workloads} predicts would move workloads between the pressure regimes---is left to future work.

\textbf{Baseline fidelity.} Hawkeye and Mockingjay are our reimplementations from the original papers; both reproduce their expected behaviour on capacity-contended workloads (e.g., Hawkeye is the best baseline on 4-core canneal). Their degradation under sharing is consistent across all workloads and scales, which argues against an implementation artefact, but independent replication would strengthen this finding.

\section{Related Work}
\label{sec:related}

\textbf{Heuristic replacement.} LRU and its approximations dominate practice; RRIP~\cite{jaleel2010rrip} re-references blocks by predicted interval and remains, with SHiP's PC-signature insertion~\cite{wu2011ship}, the strongest cheap baseline---a position our ocean results reaffirm under sharing.

\textbf{Learned replacement.} Sampling dead-block prediction~\cite{khan2010sdbp} and reuse-distance--based replacement~\cite{keramidas2007rdip} established the prediction framing that the modern line of work builds on. Hawkeye~\cite{jain2016hawkeye} casts replacement as supervised imitation of Belady's MIN~\cite{belady1966study}; Mockingjay~\cite{shah2022mockingjay} sharpens the target from binary labels to reuse distances; Glider~\cite{shi2019glider} distils offline deep models into deployable predictors; multiperspective reuse prediction~\cite{jimenez2017multiperspective} and perceptron-based reuse prediction bring the hashed-perceptron idiom~\cite{jimenez2001perceptron} to caches, which COALESCE adopts. Reinforcement learning~\cite{sethumurugan2021designing,souza2024rl} and concurrency-aware perceptrons~\cite{wu2025camp} have also been applied; of these, Souza and Freitas~\cite{souza2024rl} is closest to our setting (multicore, RL), though without coherence-state features or multithreaded-sharing evaluation. All evaluate on single-threaded or multi-programmed workloads; none observe coherence state. To our knowledge COALESCE is the first replacement policy to use MESI state and sharer count as learned features, and this paper is the first head-to-head evaluation of learned replacement under genuine multithreaded sharing.

\textbf{Coherence-aware cache management.} Coherence information has been exploited for purposes other than replacement: deactivating coherence for private blocks to shrink directories~\cite{cuesta2011coherence}, OS-assisted classification for block placement and replication in NUCA caches~\cite{hardavellas2009rnuca}, and measurement studies quantifying coherence's performance impact on real multiprocessors~\cite{molka2009memory}. These works treat coherence metadata as a management signal but leave the replacement policy untouched; COALESCE closes that gap.

\textbf{Multithreaded benchmarking.} PARSEC~\cite{bienia2008parsec} and SPLASH-3~\cite{sakalis2016splash3} (a data-race-free revision of SPLASH-2~\cite{woo1995splash2}) are the standard shared-memory suites; our Pin-based~\cite{luk2005pin} per-thread tracing and shared-VMEM replay bring them to ChampSim-class replacement studies for the first time.

\section{Conclusion}

We presented the first evaluation of learned LLC replacement under genuine multithreaded data sharing, enabled by a shared-virtual-memory overlay for ChampSim, and a new coherence-aware perceptron policy, COALESCE, designed for that regime. The evaluation yields three findings. First, the multi-programmed state of the art does not transfer to our workloads: Mockingjay places last on three of the four tested and Hawkeye falls behind LRU on half of them---evidence that on multithreaded workloads, evaluation methodology rather than model capacity is a primary limiter of learned replacement. Second, coherence features earn their cost exactly where coherence is real: COALESCE wins outright on the write-heavy irregular canneal (5\% over Hawkeye, 30\% over Mockingjay at 8 cores) and its sharer-count feature contributes a reproducible 7\% on genuinely-shared workloads while remaining inert elsewhere. Third, the boundaries are predictable: regular access patterns favour RRIP heuristics, read-only LLC traffic disables coherence features, and low LLC pressure makes all policies equivalent---a three-axis characterisation that practitioners can apply before deploying any learned policy.

Future work follows directly: multithreaded GAP and coherence-stress SPLASH workloads to populate the high-sharing corner of the characterisation; LLC size sweeps; a directory-accurate coherence model; and, most promisingly, a policy designed \emph{from} the characterisation---one that learns online which of its feature axes the running workload activates, rather than weighting them statically.

\begin{thebibliography}{20}

\bibitem{belady1966study}
L. A. Belady,
``A study of replacement algorithms for a virtual-storage computer,''
\textit{IBM Systems Journal}, vol. 5, no. 2, pp. 78--101, 1966,
doi: \href{https://doi.org/10.1147/sj.52.0078}{10.1147/sj.52.0078}.

\bibitem{jaleel2010rrip}
A. Jaleel, K. B. Theobald, S. C. Steely Jr., and J. Emer,
``High performance cache replacement using re-reference interval prediction (RRIP),''
in \textit{Proc. 37th Int. Symp. Computer Architecture (ISCA)},
Saint-Malo, France, 2010, pp. 60--71,
doi: \href{https://doi.org/10.1145/1815961.1815971}{10.1145/1815961.1815971}.

\bibitem{wu2011ship}
C.-J. Wu, A. Jaleel, W. Hasenplaugh, M. Martonosi, S. C. Steely Jr., and J. Emer,
``SHiP: Signature-based hit predictor for high performance caching,''
in \textit{Proc. 44th IEEE/ACM Int. Symp. Microarchitecture (MICRO)},
Porto Alegre, Brazil, 2011, pp. 430--441,
doi: \href{https://doi.org/10.1145/2155620.2155671}{10.1145/2155620.2155671}.

\bibitem{jain2016hawkeye}
A. Jain and C. Lin,
``Back to the future: Leveraging Belady's algorithm for improved cache replacement,''
in \textit{Proc. 43rd Int. Symp. Computer Architecture (ISCA)},
Seoul, South Korea, 2016, pp. 78--89,
doi: \href{https://doi.org/10.1109/ISCA.2016.17}{10.1109/ISCA.2016.17}.

\bibitem{shah2022mockingjay}
I. Shah, A. Jain, and C. Lin,
``Effective mimicry of Belady's MIN policy,''
in \textit{Proc. IEEE Int. Symp. High-Performance Computer Architecture (HPCA)},
Seoul, South Korea, 2022, pp. 558--572,
doi: \href{https://doi.org/10.1109/HPCA53966.2022.00048}{10.1109/HPCA53966.2022.00048}.

\bibitem{khan2010sdbp}
S. M. Khan, Y. Tian, and D. A. Jiménez,
``Sampling dead block prediction for last-level caches,''
in \textit{Proc. 43rd IEEE/ACM Int. Symp. Microarchitecture (MICRO)},
Atlanta, GA, USA, 2010, pp. 175--186,
doi: \href{https://doi.org/10.1109/MICRO.2010.24}{10.1109/MICRO.2010.24}.

\bibitem{keramidas2007rdip}
G. Keramidas, P. Petoumenos, and S. Kaxiras,
``Cache replacement based on reuse-distance prediction,''
in \textit{Proc. 25th Int. Conf. Computer Design (ICCD)},
Lake Tahoe, CA, USA, 2007, pp. 245--250,
doi: \href{https://doi.org/10.1109/ICCD.2007.4601909}{10.1109/ICCD.2007.4601909}.

\bibitem{shi2019glider}
Z. Shi, X. Huang, A. Jain, and C. Lin,
``Applying deep learning to the cache replacement problem,''
in \textit{Proc. 52nd IEEE/ACM Int. Symp. Microarchitecture (MICRO)},
Columbus, OH, USA, 2019, pp. 413--425,
doi: \href{https://doi.org/10.1145/3352460.3358319}{10.1145/3352460.3358319}.

\bibitem{jimenez2001perceptron}
D. A. Jiménez and C. Lin,
``Dynamic branch prediction with perceptrons,''
in \textit{Proc. 7th IEEE Int. Symp. High-Performance Computer Architecture (HPCA)},
Monterrey, Mexico, 2001, pp. 197--206,
doi: \href{https://doi.org/10.1109/HPCA.2001.903263}{10.1109/HPCA.2001.903263}.

\bibitem{jimenez2017multiperspective}
D. A. Jiménez,
``Multiperspective reuse prediction,''
in \textit{Proc. 50th IEEE/ACM Int. Symp. Microarchitecture (MICRO)},
Cambridge, MA, USA, 2017, pp. 436--448,
doi: \href{https://doi.org/10.1145/3123939.3123942}{10.1145/3123939.3123942}.

\bibitem{sethumurugan2021designing}
S. Sethumurugan, J. Yin, and J. Sartori,
``Designing a cost-effective cache replacement policy using machine learning,''
in \textit{Proc. IEEE Int. Symp. High-Performance Computer Architecture (HPCA)},
Seoul, South Korea, 2021, pp. 291--303,
doi: \href{https://doi.org/10.1109/HPCA51647.2021.00033}{10.1109/HPCA51647.2021.00033}.

\bibitem{souza2024rl}
M. A. Z. Souza and H. C. Freitas,
``Reinforcement learning-based cache replacement policies for multicore processors,''
\textit{IEEE Access}, vol. 12, 2024.

\bibitem{wu2025camp}
Y. Wu \textit{et al.},
``Concurrency-aware cache miss cost prediction with perceptron learning,''
in \textit{Proc. ACM Great Lakes Symp. VLSI (GLSVLSI)}, 2025.

\bibitem{cuesta2011coherence}
B. Cuesta, A. Ros, M. E. Gómez, A. Robles, and J. Duato,
``Increasing the effectiveness of directory caches by deactivating coherence for private memory blocks,''
in \textit{Proc. 38th Int. Symp. Computer Architecture (ISCA)},
San Jose, CA, USA, 2011, pp. 93--104,
doi: \href{https://doi.org/10.1145/2000064.2000076}{10.1145/2000064.2000076}.

\bibitem{hardavellas2009rnuca}
N. Hardavellas, M. Ferdman, B. Falsafi, and A. Ailamaki,
``Reactive NUCA: Near-optimal block placement and replication in distributed caches,''
in \textit{Proc. 36th Int. Symp. Computer Architecture (ISCA)},
Austin, TX, USA, 2009, pp. 184--195,
doi: \href{https://doi.org/10.1145/1555754.1555779}{10.1145/1555754.1555779}.

\bibitem{molka2009memory}
D. Molka, D. Hackenberg, R. Schöne, and M. S. Müller,
``Memory performance and cache coherency effects on an Intel Nehalem multiprocessor system,''
in \textit{Proc. 18th Int. Conf. Parallel Architectures and Compilation Techniques (PACT)},
Raleigh, NC, USA, 2009, pp. 261--270,
doi: \href{https://doi.org/10.1109/PACT.2009.22}{10.1109/PACT.2009.22}.

\bibitem{sorin2011coherence}
D. J. Sorin, M. D. Hill, and D. A. Wood,
\textit{A Primer on Memory Consistency and Cache Coherence},
Synthesis Lectures on Computer Architecture, Morgan \& Claypool, 2011.

\bibitem{hennessy2017quantitative}
J. L. Hennessy and D. A. Patterson,
\textit{Computer Architecture: A Quantitative Approach}, 6th ed.,
Morgan Kaufmann, 2017.

\bibitem{bienia2008parsec}
C. Bienia, S. Kumar, J. P. Singh, and K. Li,
``The PARSEC benchmark suite: Characterization and architectural implications,''
in \textit{Proc. 17th Int. Conf. Parallel Architectures and Compilation Techniques (PACT)},
Toronto, Canada, 2008, pp. 72--81,
doi: \href{https://doi.org/10.1145/1454115.1454128}{10.1145/1454115.1454128}.

\bibitem{sakalis2016splash3}
C. Sakalis, C. Leonardsson, S. Kaxiras, and A. Ros,
``Splash-3: A properly synchronized benchmark suite for contemporary research,''
in \textit{Proc. IEEE Int. Symp. Performance Analysis of Systems and Software (ISPASS)},
Uppsala, Sweden, 2016, pp. 101--111,
doi: \href{https://doi.org/10.1109/ISPASS.2016.7482078}{10.1109/ISPASS.2016.7482078}.

\bibitem{woo1995splash2}
S. C. Woo, M. Ohara, E. Torrie, J. P. Singh, and A. Gupta,
``The SPLASH-2 programs: Characterization and methodological considerations,''
in \textit{Proc. 22nd Int. Symp. Computer Architecture (ISCA)},
Santa Margherita Ligure, Italy, 1995, pp. 24--36,
doi: \href{https://doi.org/10.1145/223982.223990}{10.1145/223982.223990}.

\bibitem{luk2005pin}
C.-K. Luk \textit{et al.},
``Pin: Building customized program analysis tools with dynamic instrumentation,''
in \textit{Proc. ACM SIGPLAN Conf. Programming Language Design and Implementation (PLDI)},
Chicago, IL, USA, 2005, pp. 190--200,
doi: \href{https://doi.org/10.1145/1065010.1065034}{10.1145/1065010.1065034}.

\bibitem{champsim2022}
N. Gober, G. Chacon, L. Wang, P. V. Gratz, D. A. Jiménez, E. Teran, S. Pugsley, and J. Kim,
``The Championship Simulator: Architectural simulation for education and competition,''
\textit{arXiv preprint arXiv:2210.14324}, 2022,
doi: \href{https://doi.org/10.48550/arXiv.2210.14324}{10.48550/arXiv.2210.14324}.

\end{thebibliography}

\end{document}
